\subsection{The algebraic stack comparison in every stable type}

\begin{result}{E.T4}{Proved}{The structured geometric moduli automorphism groups}
Let $A_{g,n}$ be the group of $\overline\Q$-linear algebraic
autoequivalences of the stack of smooth ordered $n$-pointed
genus-$g$ curves, up to $2$-isomorphism.  For every stable type,
\[
\begin{array}{c|c|c}
 (g,n)&A_{g,n}&\ker(S_n\longrightarrow A_{g,n})\\\hline
 (0,3)&1&S_3\\
 (0,4)&S_3&V_4\\
 (1,1)&C_2&1\\
 (1,2)&1&S_2\\
 (2,0)&C_2&1\\
 \text{every other stable type}&S_n&1
\end{array}
 \tag{E.T5}
\]
The two displayed $C_2$ groups are generated by the explicit
quadratic twists described below.  Thus, under this structured
interpretation, $G_{g,n}=\operatorname{Im}S_n$ holds in every
stable type except $(1,1)$ and $(2,0)$.  The proposed stable-type
list of failures of faithfulness also needs $(1,2)$.
The forgetful map from this group to the unringed topos group is
the first map in (E.T1) at $(0,4)$; it retains the algebraic
source of each automorphism and does not exhaust that target.
\end{result}

\begin{proof}
We first give the analytic upper bound, with the exact maps
that transfer it to algebraic stacks.  The Teichm\"uller
space $\mathcal T_{g,n}$ is the simply connected analytic
orbifold cover of $\mathcal M_{g,n}^{\rm an}$, with deck group
the pure mapping class group.  An analytic stack equivalence
lifts to a biholomorphic equivalence of this cover: pull back
the universal covering, choose one lift of a basepoint, and
use simple connectedness.  The local etale charts make the
lift holomorphic.  Two choices differ by a deck transformation.
For a faithful deck action this gives exactly
\[
 \Aut(\mathcal M_{g,n}^{\rm an})
   =N_{\Aut_{\rm hol}(\mathcal T_{g,n})}(\Gamma_{g,n})
      /\Gamma_{g,n}.
 \tag{E.T6}
\]
An equivariant lift conversely descends, so this formula has
both directions.

The equality of Teichm\"uller and Kobayashi metrics is
\cite[Theorem 4, p.~936]{D_EarleKraKrushkal1994}; its proof
from the disk-lifting Theorem~2 and Proposition~1 was opened.
Biholomorphisms preserve Kobayashi distance in both directions.
Ivanov's Theorem~A, including its exact exceptional types,
therefore gives the full biholomorphism group as the
orientation-preserving full mapping class group in the
nonexceptional cases \cite[Theorem A and \S5]{Ivanov2001Royden}.
Orientation-reversing mapping classes act antiholomorphically;
on a positive-dimensional Teichm\"uller space they cannot also
be holomorphic.  The pure group is normal in the full group,
with quotient $S_n$.  Its action is faithful outside the
listed low types; a mapping class in the kernel represents
an automorphism of every marked curve, and the generic
curve has none in those types.  Formula (E.T6) gives $S_n$.
Each element is realized algebraically by relabelling.

Analytification does not introduce an unnoticed kernel in
this deduction.  For two algebraic equivalences, their
isomorphism functor over an etale atlas is finite: the
relative isomorphism scheme of stable smooth pointed curves
is finite.  A section of its analytification is algebraic.
On the finite etale locus this is the finite-cover comparison;
its graph then extends uniquely over the normal atlas by
finite birational normalization.  The sections agree on
overlaps and descend.  Thus an analytic $2$-isomorphism here
comes from an algebraic one.  This proves the algebraic upper
bound; the exhibited algebraic relabellings prove equality.

The two genus-zero exceptions were proved in E.T2, with
$\mathcal M_{0,3}=\operatorname{Spec}k$.  We treat the other
exceptions with their stabilizers retained.

\emph{Two marks in genus one.}
The canonical involution of $(E,p,q)$ given in group-law
coordinates by $x\mapsto p+q-x$ exchanges the two marks.
Its quotient has four branch points and the distinct fifth
point which is the common image of $p,q$.  This gives an
equivalence
\[
       \mathcal M_{1,2}\simeq[\mathcal M_{0,5}/S_4].
 \tag{E.T7}
\]
Here the fifth point is fixed.  The inverse in families has
an explicit square-root construction.  Write
$\pi:P\to S$ for the genus-zero family, $D$ for its four-point
branch divisor, and $r:S\to P$ for the disjoint fifth section.
Put
\[
 N=r^*\mathcal O_P(r),\qquad
 L=\mathcal O_P(2r)\otimes\pi^*N^{-2}.
\]
The degree-zero line bundle
$\mathcal O_P(D)\otimes L^{-2}$ comes from $S$: this follows
on a local trivialization of the $\mathbf P^1$-bundle from
$\operatorname{Pic}(\mathbf P^1_S)/\operatorname{Pic}(S)=\Z$,
and the isomorphisms descend by evaluation at $r$.
Both factors have their canonical trivialization along $r$,
so this evaluation identifies $L^2$ with $\mathcal O_P(D)$.
The section defining $D$ then has value $1$ along $r$.
The double cover
\[
 \operatorname{Spec}_P(\mathcal O_P\oplus L^{-1}),
 \qquad u^2=s_D,
\]
has the ordered sections $u=1,-1$ above $r$.  Its relative
genus is one by the degree-four branch formula.  Choosing
the other local square root interchanges the sections by
the deck involution, giving the unique isomorphism compatible
with the chosen first section.  These isomorphisms obey the
cocycle identity.  The construction inverts the quotient
functor on objects and isomorphisms, including over bases
where a local square root was initially required.
The deck subgroup in the Teichm\"uller cover of (E.T7) is the
inverse image of $S_4<S_5$ in the full five-puncture group.
Theorem~A applies to this sphere.  Its normalizer quotient is
$N_{S_5}(S_4)/S_4=1$: a permutation normalizing a point
stabilizer fixes its uniquely fixed point.  Thus the entire
automorphism group is trivial.  The displayed involution
$x\mapsto p+q-x$ is also the natural $2$-isomorphism making
the mark transposition trivial.

\emph{One mark in genus one.}
Write an elliptic object in short Weierstrass form as
\[
 (L,a,b),\qquad a\in L^4,\quad b\in L^6,
 \quad\Delta=4a^3+27b^2\in L^{12}\text{ nowhere zero}.
\]
This is the quotient presentation with weights $4,6$.
Rigidifying the generic $\mu_2$ gives
\[
 \mathcal B_1=
 [\operatorname{Spec}k[a,b,\Delta^{-1}]/\mathbf G_m],
       \qquad\text{weights }2,3.
\]
Its coarse space is the $j$-line, with distinct stabilizer
orders three and two at $0$ and $1728$.  A coarse
automorphism fixes these two points and hence is the
identity affine transformation.  Over the identity coarse
map an automorphism of this effective root orbifold is
$2$-isomorphic to the identity: in a local root chart the
faithful tangent character forces the identity map on
the cyclic stabilizer.  The local isomorphisms are unique
on the dense locus with trivial inertia and therefore agree
on overlaps.  Thus $\Aut(\mathcal B_1)=1$.

For a $\mu_2$-banded gerbe over an effective base, the
autoequivalences over the identity base are
$H^1(\mathcal B,\mu_2)$.  To see this explicitly, choose
local objects of the gerbe.  The images under an
equivalence differ from these objects by $\mu_2$-torsors.
Their transition isomorphisms form a torsor on the base;
changing the equivalence by a $2$-isomorphism changes this
cocycle by a coboundary.  Conversely contracting each local
object with a base torsor gives an equivalence, and these
constructions are inverse.  There is no extra band
automorphism since $\Aut(\mu_2)=1$.

The coordinate ring $k[a,b,\Delta^{-1}]$ is a unique
factorization domain and its units are $k^*\Delta^{\Z}$.
Equivariant line bundles are therefore characters of
$\mathbf G_m$, modulo the weight six of $\Delta$; hence
$\operatorname{Pic}(\mathcal B_1)=\Z/6$.
The units on the quotient are $k^*$, all squares.
The Kummer sequence gives
$H^1(\mathcal B_1,\mu_2)=\operatorname{Pic}(\mathcal B_1)[2]
\simeq C_2$.  Every automorphism of $\mathcal M_{1,1}$
descends to the rigidification, so this computes its full
group.  An explicit generator is
\[
       (L,a,b)\longmapsto
       (L^7,\Delta^2a,\Delta^3b).
 \tag{E.T8}
\]
Its discriminant is $\Delta^7$; its square is
$(L^{49},\Delta^{16}a,\Delta^{24}b)$, canonically identified
with the starting object by multiplication by $\Delta^{-4}$.
On inertia of order four its map is $u\mapsto u^7=u^{-1}$,
whereas on inertia of order six it is the identity.
Conjugation in these cyclic stabilizers is trivial, so the
order-four action proves that (E.T8) is not $2$-isomorphic
to the identity.

\emph{The closed genus-two stack.}
Every genus-two curve is hyperelliptic.  Quotient by its
canonical involution identifies its rigidification with
\[
          \mathcal B_2=[\mathcal M_{0,6}/S_6].
\]
The full sphere mapping class group is the deck group for
its Teichm\"uller cover and is also the full biholomorphism
group of that cover.  Thus (E.T6) gives
$\Aut(\mathcal B_2)=1$.  The map
$\mathcal M_2\to\mathcal B_2$ is a $\mu_2$-banded gerbe:
locally a six-point branch divisor has a double cover,
and its automorphisms over the branch object are exactly
the two deck maps.  The preceding cocycle calculation
reduces its automorphism group to
$H^1(\mathcal B_2,\mu_2)$.

The full sphere six-point group has the braid generators
$\sigma_1,\ldots,\sigma_5$, the sphere relation of exponent
sum $10$, and the full-twist relation of exponent sum $30$.
The braid relations identify all generators in its
abelianization.  Thus its abelianization is exactly
$\Z/\gcd(10,30)=\Z/10$, and its characters to $C_2$ form
$C_2$.  The nontrivial character is the sign of the six-point
permutation.  Riemann existence identifies these characters
with $H^1(\mathcal B_2,\mu_2)$.
The generator twists the hyperelliptic double cover by the
sign torsor of its branch divisor; contracting by that
torsor twice is canonically the identity.  The torsor is
nontrivial, since the connected cover
$[\mathcal M_{0,6}/A_6]\to\mathcal B_2$ has degree two.
The cocycle correspondence proves that its gerbe
autoequivalence is nontrivial.  It is defined algebraically
over $\Q$.  This finishes the two remaining entries and
all the maps from $S_n$ in the table.
\end{proof}

\begin{result}{E.T5}{Proved}{The three unstable genus-zero algebraic objects}
For smooth proper genus-zero curves with zero, one or two
ordered marks, the algebraic moduli stacks over
$\overline\Q$ and their groups of automorphism classes are
\[
\begin{array}{c|c|c}
 (0,0)&B\operatorname{PGL}_2&1\\
 (0,1)&B(\mathbf G_a\rtimes\mathbf G_m)&1\\
 (0,2)&B\mathbf G_m&C_2.
\end{array}
\]
In the last row the mark transposition gives inversion
of $\mathbf G_m$ and is faithful.  These statements use
the displayed algebraic classifying stacks; their
connected stabilizers and the corresponding lisse-etale
objects are retained.
\end{result}

\begin{proof}
Etale locally, trivialize the genus-zero curve and send its
ordered marks to the usual fixed points of $\mathbf P^1$.
Its automorphism group is successively
$\operatorname{PGL}_2$, the affine group, and $\mathbf G_m$,
which proves the three classifying-stack identifications.
For a group $H$ over the algebraically closed field $k$,
an equivalence of $BH$ sends the trivial torsor to a
torsor over $k$, which is trivial here.  Choosing its
trivialization identifies the equivalence with an
automorphism of $H$; changing the trivialization is inner
conjugation.  Thus its group is $\Out_k(H)$.

The group $\operatorname{PGL}_2$ has only inner algebraic
automorphisms: its automorphism permutes its Borel groups,
giving a transformation of their parameter scheme
$\mathbf P^1$; conjugate by that fractional-linear
transformation, then the automorphism fixes the two root
groups and torus, whose conjugation relations force the
remaining root scalings to come from that torus.
For the affine group its characteristic unipotent radical
is $\mathbf G_a$.  An automorphism scales this radical by
some $c\in k^*$; the conjugation character forces the
induced torus map to be the identity.  The remaining
torus-to-unipotent cocycle is $u\mapsto b(1-u)$, as follows
by applying the cocycle identity in the two orders to
$u,v$.  Scaling and this cocycle are respectively
conjugation by a dilation and a translation.  Its outer
group is therefore trivial.  Finally the algebraic
automorphisms of $\mathbf G_m$ are $z\mapsto z$ and
$z\mapsto z^{-1}$, and inner conjugation is trivial.
Interchanging $0$ and $\infty$ conjugates the dilation
action to its inverse, proving the asserted map of
mark permutations.
\end{proof}
